\pagebreak

\chapter{Program Source Structure}\label{chap:progsrcstruct}
    
    Module, procedure, function and class definitions are processed during
    parsing, not during immediate execution.

    All statements are executed by the 3PL interpreter. Statements which when
    executed produce logic for the FPGA are called {\bf target} statements. All
    other statements are called {\bf immediate} statements.

    Statements are refered to as {\bf top level} if they are
    not enclosed in a control structure (if(), for(), while() etc.).

    Top level immediate statements are executed in the order in which they are
    encountered. Top level statements may of course be control structures,
    inline modules or module or procedure calls in which case nested
    statements will be executed.
    
    Module, procedure and function definitions may be distributed amongst
    top-level statements (but not at a deeper level of nesting) since they have
    been removed by the parser. Code is probably more readable though if
    definitions are not followed by statements.
    

    \section{scope}

        \index{scope}
        \index{variable!scope}
        The {\bf 3PL} compiler is given a source file name via the command line.
        The file is treated as a {\bf file module}, by which is meant that when
        code from that file is executed it is treated as though a module had
        been called, the module name being the file name. A file module has the
        property that it provides a static scope which is accessible to any code
        within the file, including within any inline modules or within
        procedures, functions or modules defined in the file module. That scope
        is called the {\bf file scope}.

        If a 3PL source file inserts another file by means of a {\bf source}
        statement then that inserted file is expanded in place.

        If a 3PL source file inserts another file by means of a {\bf module}
        statement then that file is also expanded in place but it is treated
        as a file module, that is it behaves as though it has been called as a
        module at that point in the code, in this case a file module.

        Nesting of {\bf source} and {\bf module} statements can be to any
        depth (but note that a given file can only be inserted once via a
        source or module statement (if source code is required in more than one
        place it should be in a named module or procedure).

        File scope has the property that it is accessible to any code within
        the file (including code inserted by {\bf source} statements but
        not to nested file module statements). Every statement has an associated
        static file scope which is that of the file module within which it is
        contained.
        
        If a class is created then its code body has a class scope and a common
        class scope. These are accessible within the class body and any modules,
        procedures and functions defined in that block. It is not accessible to
        any nested class blocks.

        Finally there is a global scope and variables can be explicitly created
        in that global scope. Declared modules, procedures and functions can
        also be explicitly placed in global scope. Directives are only placed
        in global scope.

        Scopes can be one of the following -
        \blist
        \item   global - accessible anywhere
        \item   file - the scope of a module created by inserting a file, the
                main source file or a file imported via a module statement,
                and accessible anywhere within it (including within any
                defined modules, procedures or functions, but not within
                deeper nested file modules).
        \item   class - the scope of the current class block if applicable
        \item   local - the scope of a procedure, function or module.
        \item   block - the scope of a code block in a control structure.
        \item   calling - the scope of the calling module, procedure or
                function.
        \blend
        
        File scope is the scope of a file module, a module associated with the
        main source file or imported by a module statement as described
        previously. This is a static scope that can be accessed by any code
        within the file module, including code within modules, procedures or
        functions defined within that file module regardless of from where they
        are called.
        
        Where file modules are nested, code within a nested scope can only
        access the nested file module scope, not that of enclosing file modules.
        
        Global scope, file scope are static scopes, i.e. they are
        established during compilation and are not necessarily on the call stack.
        
        Class scopes are created when class instances are created and are
        accessible during execution of a class code body or when a class
        variable or class module, procedure or function is called.

        When a variable is referenced or evaluated the identifier is searched
        for through scopes in the following order (where they exist) -
        \blist
        \item   nested block scopes back up to the local scope
        \item   local scope (module, procedure or function) - the local scope
                will be the file scope if we are in a file module
        \item   the current class scope if applicable
        \item   the current file scope
        \item   global scope
        \blend
        The search completes when the first instance of the variable
        identifier is encountered.
        
        In some cases the local scope will be the file scope. In-line statements
        in a file module will have a file module scope which is the same as
        the local scope for the module.

        Note that inline modules are an exception to this rule. The scope
        search continues past the local scope of the module back to the scope
        of the containing named module, procedure or function. This gives
        in-line module statements access to the enclosing scope.

        When a variable is created, by default it is created in the current
        scope. If within a code block in braces, the variables will be created
        in the scope of that block. Thus to conditionally create variables
        in the LOCAL scope they must be in separate conditional statements, e.g. -
        \begin{lstlisting}[gobble=8, language=threepl]
            if (l)
                int(i);
            if (l)
                str(s);
        \end{lstlisting}
        or alternatively the calling scope can be specified -
        \begin{lstlisting}[gobble=8, language=threepl]
            if (l) {
                int(calling.i);
                str(calling.s);
            }
        \end{lstlisting}
        Note that if "class" is specified as the scope this only accesses the
        class instance scope, not the class common scope.

        The previous example shows that variables can be created in a
        specified scope by prepending a scope specifier as follows -
        \blist
        \item   local. - use the scope of the current module, procedure or
                function. (this may be file scope)
        \item   calling. - use the scope of the calling module, procedure or
                function.
        \item   class. - use the current class scope if applicable.
        \item   file. - use the current static file module scope.
        \item   global. - use the static global scope.
        \blend

        Similarly when evaluating variables the scope can be explicitly
        specified (for example to disambiguate variables with the same
        identifier but in different scopes, or in the case of calling scope to
        access variables not normally accessible).

        As an example the source file s1.3pl contains -
        \begin{lstlisting}[gobble=8, language=threepl]
           int(s);
           int(global.g1);
           module "s2.3pl"

           p();
           [
               int(j);
           ]

           procedure p () {
               int(l);
               if (true) {
                   int(b);
               }
           }

           global procedure q() {
               .
               .
           }
        \end{lstlisting}
        and source file s2.3pl -
        \begin{lstlisting}[gobble=8, language=threepl]
           int(s);
           int(global.g2);
        \end{lstlisting}
        Execution begins in s1.3pl. It creates an integer {\bf s} which is in
        the file module s1 scope, then an integer {\bf g1} which is explicitly
        placed in global scope. Execution now enters source file s2.3pl which is
        a new file module. An integer {\bf s} which is in file module s2 scope
        and an integer {\bf g2} which is in global scope are created. Integers
        {\bf g1} and {\bf g2} can be accessed in either s1.3pl or s2.3pl as they
        are global, but each file module scope now contains a unique variable
        {\bf s} which can only be accessed from within that file module.

        We have reached the end of s2.3pl and now continue execution in s1.3pl.
        Integer {\bf s} in file module s2 is now no longer accessible. Procedure
        {\bf p} is called. This procedure is declared within the s1 file module
        so could only be called from within that file module, but note that
        procedure {\bf q} could be called from anywhere as it has been declared
        global. Within procedure {\bf p} integer {\bf l} is created and that is
        in local scope, then integer {\bf b} is created in the scope of the {\em
        if} block. Within the {\em if} block variables {\bf g1}, {\bf g2}, {\bf
        l} and {\bf b} are accessible and variable {\bf s} in file module s1 is
        also accessible. On leaving the {\em if} block variable {\bf b} would no
        longer be accessible (though we could still use it if we had a pointer
        to it). After returning from procedure {\bf p} we enter an inline module
        in which  integer {\bf j} is created. {\bf j} is local to the inline
        module and at this point {\bf g1} and {\bf g2} can still be accessed, as
        can integer {\bf s} in file module s1.

        The use of {\bf source} and {\bf module} statements and scope specifiers
        such as {\bf global.} and {\bf file.} provide a mechanism to achieve
        similar scope control to that provided by {\bf static} and {\bf extern}
        in {\bf C}. The need for a different mechanism arises from the fact that
        variables in 3PL are created during immediate execution, not declared
        during parsing. For the same reason we cannot simply have a list of
        source files as in {\bf C} but must use {\bf source} and {\bf module}
        statements to link source files together in such a way that variable
        creation occurs in the correct order.
        
        File scope is intended to be used to hide variables such as I/O pins and
        signals or intermediate clock interconnections. All I/O pin locations
        strings, clocks, external memory interfaces and external bus interfaces can
        be coded in a source file included using a module statement. File scope
        variables within that file will only be accessible from code within the
        file. Any variables that need to be accessible from elsewhere can be made
        global. The interface and other services provided by code within the file
        can be made available through module, procedure and function calls which
        are declared global.
        
        Note that block scope only applies to control structures which have
        code blocks enclosed within braces ({\em \{\dots\}}.
        Conditional variable creation can be implemented by the use
        of an {\em if} containing a single statement, e.g.
        \begin{lstlisting}[gobble=8, language=threepl]
           if (condition)
               int(i);
        \end{lstlisting}
        which will create the integer variable {\em i} in the scope
        containing the {\em if} if logical variable {\em condition} is true.
        As shown in an earlier example it would be a mistake to write
        \begin{lstlisting}[gobble=8, language=threepl]
           if (condition) {
               int(i);
           }
        \end{lstlisting}
        as the variable would be created in the scope within the {\em if} and
        would vanish when the {\em if} completed. This should be written -
        \begin{lstlisting}[gobble=8, language=threepl]
           if (condition) {
               int(local.i);
           }
        \end{lstlisting}
        for the scope of the enclosing module, procedure or function or
        \begin{lstlisting}[gobble=8, language=threepl]
           if (condition) {
               int(calling.i);
           }
        \end{lstlisting}
        for the scope immediately enclosing the {\bf if()}. Clearly this is
        the way to conditionally create a group of variables where more than
        one variable creation call is needed. 
        
        If a procedure is to be written to create a variable in the current scope
        then the variable must be created in the calling scope, e.g.
        \begin{lstlisting}[gobble=8, language=threepl]
               procedure
               my_var (str(name)) {
                   name = "calling."+name;
                   int(name->);
               }
        \end{lstlisting}
        Here to make the example non-trivial we have passed the identifier of an
        integer we wish to create in the current (i.e. calling) scope. The
        procedure prepends the variable name string with keyword "calling." and
        then creates an integer using the indirect operator on the string (an
        indirect operator is normally used on a pointer variable but it can also
        be used on a string variable). Note that the procedure parameter is
        immediate mode so the argument will be passed by value, i.e. copied to
        variable {\em name} which is in effect a local variable. Thus we can
        change the value of the parameter within the procedure body.

        \index{nested declarations}
        Procedures, functions and named modules can be declared within any
        module, procedure or function. By default they can only be called within
        the same module, procedure or function or from within any module,
        procedure or function called from the current scope. A procedure,
        function or named module declared within a file module can be made
        globally accessible by preceding the {\bf module}, {\bf procedure} or
        {\bf function} keyword by the keyword {\bf global}. This cannot be done
        when they are declared at a deeper level.

        Variables in various scopes can be printed by library procedure {\bf
        printvars()} \autorefph{sec:lpprintvars} in {\bf stdlib.3pl}. Among other
        parameters this lists the scopes in which variables occur. The scopes can
        be selected by an argument to the procedure.

        Directives can be printed by library procedure {\bf printdirs()}
        \autorefph{sec:lpprintdirs} in {\bf stdlib.3pl}.


        
    \section{immediate execution}
    
        3PL is parsed to produce a code tree for immediate execution.
        Any {\bf source} statements encountered are expanded in place.
        Any {\bf module} statements encountered are also expanded in place
        but as a form of in-line module called a {\bf file module} where the
        file contents form the body of the module.
        
        During parsing no variables are defined but module, procedure and
        function definitions are processed and entered into maps in the
        enclosing scopes in which the definitions are encountered.
        
        During immediate execution statements are executed starting at the
        beginning of the main source file. {\bf source} statements are no
        longer recognised as they have been expanded in place and are executed
        as though part of the enclosing source code. {\bf module} statements
        and in-line modules are executed as though called at the point in
        the source code at which they occur. These implicit calls are without
        any parameters or arguments.
        
        All variables are created during immediate execution.
    
        \index{execution!immediate}
        Immediate execution begins with execution of the initial file module.
        Files included using {\bf source} or {\bf module} statements are
        executed at the point of inclusion.

        During immediate execution modules, procedures and functions may be
        called. For immediate execution there is no distinction between
        modules and procedures, the calling syntax and semantics being the
        same. Both follow the same scope rules. The only distinction between
        modules and procedures is that modules initiate target code execution
        streams and control the way multiple readers of a queue are handled.

        At the end of application code execution it is usual to have many
        housekeeping tasks to complete. For example a resource such as RAM or
        an external interface may have to be shared amongst multiple accesses.
        The logic to implement this can usually only be generated at the end
        of the application code. Similarly there are often variables created
        in different included files which ideally should be created very early
        in execution rather than when encountered during sequential immediate
        execution. The above can be achieved by declaring module definitions
        using keywords {\bf leadermodule} or {\bf trailermodule} (see
        \autorefph{sec:mpfcdecl}). The declaration includes a mandatory comment
        string and an optional ordering integer. Modules thus declared do not
        have an identifier or parameters as they are not called by name and
        cannot be called directly from source code. More than one leadermodule
        or trailermodule can be defined and they are placed in a list when
        encountered during parsing. When immediate execution begins the leader
        modules are executed in the order determined by the ordering integer
        or in which they were queued for the same ordering integer, the source
        code is then executed in sequential order, and finally the trailer
        modules are executed in the same way. Leader and trailer modules may
        call other modules, procedures or functions as usual. The leader or
        trailer module comment strings are printed to the report file when
        these modules are executed and this aids verification that they are
        executed in the desired order.

        It is usual for an application program to first include a header file
        module which will define a trailer module. Housekeeping variables can
        thus be hidden in the scope of the header file module to be used later
        by the trailer module. The header file module usually defines FPGA
        part number and package pins and sets up clocks. Thus an application
        program is usually of the form
        \begin{lstlisting}[gobble=8, language=threepl]
           module "header.3pl"
           .
           .
           .
        \end{lstlisting}
        where header.3pl sets up the environment for the application code
        and defines a trailer module to be called automatically at the
        completion of execution of the application code to tie up loose ends.

    \section{Target code execution}
    
        \index{execution!target}
        3PL is designed to process continuous data streams and the target
        execution control mechanism has been designed with this in mind. In
        particular it is convenient to be able to write assignment statements
        so that they are data driven. If we have a number of queue variables
        q1, q2 etc. it would be convenient to be able to write statements of
        the form
        \begin{lstlisting}[gobble=8, language=threepl]
               q1 <<- q2 + 5;
               q4 <<- q1 - q3;
        \end{lstlisting}
        and have these execute on every datum that comes down each queue,
        rather than have to explicitly enclose each in an infinite loop
        \begin{lstlisting}[gobble=8, language=threepl]
           par {
               seq while (true) {
                   q1 <<- q2 + 5;
               }
               seq while (true) {
                   q4 <<- q1 - q3;
               }
           }
        \end{lstlisting}
        
        \subsection{the role of the module}

            \index{module}
            Modules have two prime characteristics relevant to the target -   

            \blist
            \item   they constitute single sinks for queues
            \item   they initiate repeated execution of target code streams
            \blend
            
            The second of these is discussed in the next section.
            
            Queues are allowed multiple readers. A queue datum cannot be popped
            from the queue until all readers have read. However, we may wish
            to have queue reads in several places, some of them conditional.
            
            To simplify the rules for queue reading the module has been
            devised as a sink for queues.
            
            A module may contain any number of places at which a queue is read.
            Any queue read within a module receives an acknowledgement of the
            read from that module. If a module executes two or more reads of a
            queue simultaneously the queue still receives only one acknowledgement
            and all reads extract the same datum. Queue reads within a module
            which are not simultaneous extract different data.

            Where a queue is read in more than one module, each module must
            execute a read of a particular datum before it is popped and the
            next datum (if any) becomes available. Thus if the read within one
            module is delayed, other modules which have read the datum will now
            see that the queue is not available. When all modules have read the
            queue it may once again appear available to all modules. This ensures
            that all modules that contain a read operation for a queue get each
            datum regardless of when they execute the read. Note however that as
            described in the previous paragraph multiple reads of a queue within
            a module are not constrained in this way.
            
            The behaviour described above for reads within multiple modules can
            be disabled by the {\bf ignoremodules} queue variable attribute. If
            this is set to {\bf true} then all reads of the queue will be treated
            as though within a single module. The purpose of this option is to
            allow code within multiple modules to read queues while looking after
            exclusive access where necessary by means of explicit user code.
            This can also be done also by ensuring that all the queue reads are
            contained within one module but this restriction can be inconvenient
            and may result in more obscure code.
            
            
            A module is recognised as a reader for a queue if that module
            contains one or more value references for that queue. There are three
            cases which should be recognised as exceptions -
            \index{queue!availability}
            \blist
            \item   queue examines via the {\em ?} operator
            \item   queue availability tests using the {\em \verb'<'}
                    or {\em \verb'>'} operators
            \item   passing a queue to a module as an argument
            \blend
            none of these make the associated module a queue sink.
            
            A queue value reference is simply the appearance of a queue variable
            in any expression that is evaluated to read a queue datum. This might
            be the RHS of an assignment statement, a variable target array
            subscript or part of an expression tested for conditional execution.

            A queue is written by a queue assignment statement, the queue
            identifier appearing on the LHS. Each assignment to a queue
            inserts one datum into its queue. If the datum is compound (an
            array or struct) then the whole type will be written when any
            assignment is made. Multiple array members or struct fields can
            be written by separate assignment statements as long as these are
            executed simultaneously. If two assignments to the same array
            member or field are made simultaneously that is an execution
            error\footnote{at the moment target execution errors are not
            detected, but this facility may be added as a compiler option
            later}.

            Queue reads and writes from multiple parallel statements can be
            synchronised by the use of a {\bf sync \{\dots\}} block described
            in section \autorefph{sec:sync}.
            
            Unsynchronised queue writes may be arbitrated by using the {\bf
            waitpri ()} control structure, section \autorefph{sec:waitpri}
    
        \subsection{target execution streams}

            \index{execution!target!streams}
            \index{procedure}
            Procedures do not exist as far as the target is concerned.
            Procedure calls are done during immediate execution and leave no
            trace in the target code. We can conceptually eliminate the
            influence of procedure calls on target code by regarding
            procedure definitions as being expanded in place of the call with
            parameters replaced by call arguments.

            The same applies to function calls. These are again done during
            immediate execution and any resulting target code generated is
            substituted in an expression in place of the call (a function
            can also create an independent module which drives a value
            returned by the function).

            It was mentioned above that modules -
            \blist
            \item   constitute single sources and sinks for queues
            \item   initiate repeated execution of target code streams
            \blend

            The second of these characteristics is of interest here. We can
            ignore procedure calls as discussed above and consider any target
            code within a procedure as substituted for the call.

            \index{execution!target!implicit infinite loop}
            Any top level target statement in a module (remembering
            that we have conceptually eliminated procedures) is effectively
            inside its own infinite loop, that is it executes repeatedly
            forever. This applies regardless of the statement type be it a
            static or queue assignment (value, selectvalue, pointer and
            priority assignments are immediate) a seq, sync or par block, a
            when etc.

            Note that execution of each statement within its own infinite loop is
            not the same as repeated execution of a parallel block containing those
            statements. In a parallel block each statement in the block is executed
            once, but the next repetition cannot take place until every statement
            in the block has executed. For example one of the statements in the par
            block might be quite a long seq block, or an assignment statement might
            wait quite a long time for queue availability \index{queue!availability}.
            With statements each in their own infinite loop they each repeat at
            their own maximum rates.

            Thus each top level statement within a module executes repeatedly
            forever and does so independently of other top level statements.

            If a 3PL target code stream is to execute only once in the way
            conventional programming languages dictate then a stream can be
            simply terminated with a call to inbuilt procedure {\bf stop()}.

            An execution stream can proceed through a number of control
            structures. A {\bf when} conditionally selects one of two code
            blocks according to the value of a logical expression. There are
            also target {\bf while} and {\bf do while} loops. Statements can
            be enclosed in a {\bf seq} block to execute them sequentially or
            in a {\bf par} or {\bf sync} block to execute them in parallel.
            All of the above may of course be nested in any way. Sequential
            execution causes the execution stream to pass from statement to
            statement in order. Parallel execution causes the execution
            stream to diverge into multiple streams, the block waiting until
            all streams have completed before it completes.
            
            \index{execution!target!exception termination}
            \index{unless construct}
            Execution streams can be prematurely terminated by means of the {\bf
            unless} construct. This can be used with {\bf when}, {\bf branch}, {\bf
            seq}, {\bf par}, {\bf sync}, {\bf seq while}, {\bf par while}, {\bf sync
            while}, {\bf seq do while}, {\bf par do while} and {\bf sync do while}
            contructs.

            \index{execution!target!availability}
            \index{queue!availability}
            As well as an execution stream or streams there is an
            orthogonal control mechanism - queue availability. An assignment
            statement will only execute when
            \blist
            \item   execution control reaches it in its stream
            \item   all queues that it reads or writes are available
            \blend
            If any referenced queues are not available execution of that
            stream is suspended until they are. Note that this is different
            to a target conditional statement such as a {\bf when} which will
            simply choose a true or false branch of the stream but will
            not pause of its own volition (it may wait if its test expression
            contains queue reads).
        
        
        \subsection{starting target execution streams}
        

            \index{target!execution!start}
            Execution threads start in an implied infinite loop initiated within
            a module. If no special provision is made every execution thread will
            start as soon as GWE is asserted after the configuration has been
            loaded. Every clock domain has its own start pulse initiating every
            thread in that clock domain. These start pulses are initiated at the
            same time.
            
            The start of execution threads may be delayed in the
            following ways -
            \blist
            \item   after a specified number of clock pulses
            \item   after a specified time
            \item   when a target expression is true
            \blend
            In the last case the target expression must only contain signals
            that are external inputs or are derived from FPGA elements such
            as DCM lock signals since all 3PL target variables of any mode
            with state are still frozen in their initial conditions.
            
            Specification of the start condition is done using inbuilt
            procedure start() \autorefph{sec:ipstart}).
            
            The start conditions are specified relative to one clock domain and
            the start pulse in that clock domain is resynchronised into all other
            clock domains to start them.            


    \section{Clocks}

        \index{clock}
        3PL generates synchronous logic and supports any number of separate
        clocks. There are two aspects of clocks with which a program must
        deal -
        \blist
        \item   selection of the clock currently in force
        \item   generation of clock signals
        \blend

        \subsection{selecting a clock}

            \index{clock!creation}
            A clock variable is created by
            \begin{lstlisting}[gobble=12, language=threepl]
               clock(cv, ...attributes...);
            \end{lstlisting}
            The attributes arguments are optional, but if clock frequency or
            period attributes have not been set by the end of immediate execution no
            timing constraints will be generated for logic using this clock.
            \index{timing constraint}
            \index{constraint!timing}
            Only one of frequency or period attributes should be given, the other
            being derived automatically. If the duty cycle attribute is omitted 50\% will
            be assumed. In the case of an external clock signal there are
            additional attributes.
            
            The clock attributes are usually omitted when -
            \blist
            \item   the clock variable is a module, procedure or function
                    parameter, in which case it will reference the argument
            \item   the clock variable is to be assigned from another clock
                    variable using \verb':=' assignment (see below)
            \item   the clock variable is to be an argument to a library
                    module or procedure which will set the attributes
            \blend
            
            A clock variable may be assigned another clock variable following
            which it becomes an indirect reference to that variable, for example
            \begin{lstlisting}[gobble=12, language=threepl]
               clock(cv);
               cv := c;
            \end{lstlisting}
            where {\bf cv} is now a reference to {\bf c} (note: this does not
            require the use of the pointer operator). Since this is an indirect
            reference {\bf cv} can now be used in place of {\bf c}, including
            operands to equality comparisons.           
            
            The clock signal currently in force is changed by calling the
            {\bf useclock()} procedure. For example
            \begin{lstlisting}[gobble=12, language=threepl]
               useclock(c100);
            \end{lstlisting}
            would select {\bf c100} as the current clock, pushing the
            previous current clock onto a stack. The call
            \begin{lstlisting}[gobble=12, language=threepl]
               useclock();
            \end{lstlisting}
            pops the previous clock off the stack, i.e restores the previous
            current clock. If there are no previous clocks the current
            clock becomes null. Clocks can be pushed on the stack to any depth.
            
            The clock signal currently in force can only be changed at the
            top level in a module. This precludes a clock change within the
            infinite loops enclosing execution streams. An execution stream
            must have a single clock controlling all statements within it.
            Changing the clock in a module between one top level statement
            and the next poses no problem.
        
        \subsection{clock domains of variables}\label{subsec:clockdoms}

            \index{clock domain}
            Variables do not have established clock domains when created. The
            clock domain of a variable can be established in one of three ways -
            \blist
            \item   when a variable is assigned its domain is set to the current
                    clock domain, but some modes are exceptions - see below
            \item   when a variable is finally evaluated for use in an execution
                    thread its clock domain is set to the current
                    clock domain
            \item   the clock domain of a variable may be explicitly established
                    by calling inbuilt procedure {\bf attributes()} or by specifying
                    the attribute in the variable creation procedure call
                    section \autorefph{sec:ipattributes}
            \blend

            Where the clock domain of a variable is already established the
            above three cases check the clock domain instead, causing a fatal
            error if the check fails.
            
            A value variable itself has no clock domain, but the values it
            contains may have clock domains. When used as part of an evaluated
            expression the clock domains of the values assigned to the variable are
            used. If at that time they have no clock domains their clock domains are
            set to the current clock domain. Where a value variable is of compound
            type different elements may be assigned expressions with different clock
            domains. The following code illustrates this behaviour -
            \begin{lstlisting}[gobble=12, language=threepl]
               value(v, "[2]log");
               static({s1, s2, s3, s4}, "log");
               v[0] := s1;
               v[1] := s2;

               useclock(c100);
               s3 <- v[0];

               useclock(c50);
               s4 <- v[1];
            \end{lstlisting}
            The assignments to value variable {\bf v} do not establish the clock
            domains of s1 or s2. Variable {\bf v}, being value mode, does not itself
            have a clock domain. When the assignment {\bf \verb's3 <- v[0]'} is
            encountered during immediate execution the current clock domain ({\bf
            c100}) will be established for s3 (by assignment) and s1 (by evaluation
            of the expression containing v[0]). Similarly the assignment {\bf
            \verb's4 <- v[1]'} will set {\bf s4} and {\bf s2} to clock domain {\bf
            c50}. The purpose of this example is to explain that assignment of
            expressions to a value variable does not establish the clock domain
            of variables in that expression.
            
            If an expression assigned to a value variable contains variables
            whose clock domains are already established then when the value
            variable is evaluated in a clocked context the clocks will be
            checked for identity with the current clock.
            
            A value variable which is evaluated prior to its assignment is given a
            value that has the current clock domain. Later assignment to the
            variable must be in the same clock domain.

            Selectvalue, static and priority mode variables have a single clock
            domain in which both assignment and value reference are done. These
            variables have no clock domain when created and the clock domain is
            established from the current clock domain when the variable is first
            assigned or its value is used in a clocked context, e.g. an assignment
            value. All subsequent assignments or value references must be done in
            the same (current) clock domain.
            
            Input mode variables have only an output (read) clock domain since they
            can be evaluated but can not be assigned. Input mode variables have no
            clock domain when created and the clock domain is established from the
            current clock domain when the value of the variable is first used. All
            subsequent value references must be done in the same clock domain.
            
            Output mode variables have only an input (write) clock domain since they
            cannot be evaluated, only assigned. Output mode variables have no clock
            domain when created and the clock domain is established from the current
            clock domain when the value of the variable is first assigned. All
            subsequent assignments must be done in the same clock domain.        
            Note that if an output mode variable is assigned an expression
            containing value mode variables which have not been assigned values yet,
            the interim values will be set to the current clock domain. When
            ultimately these value variables are assigned values the values must
            have the same clock domain.
 
            Queue mode variables have separate read and write clocks. Again queue
            variables have no clock domains when created. The write clock
            domain is set from the current clock domain on the first write
            (assignment) to the queue after which all future writes must be in
            the same clock domain. Similarly the read clock domain is set from
            the current clock domain on the first queue read or examine and all
            further reads or examines must be done in that clock domain. Evaluation
            of expressions containing queue availability operators in a clocked
            context (i.e. but not when assigned to a value mode variable) will
            also establish or check the associated queue clock domain (read or write).
            
            The clock domain(s) of a variable can be forced by setting
            attributes {\bf domain}, {\bf readdomain} or {\bf writedomain}
            (see \autorefph{sec:attributes}) as appropriate. This pre-empts
            clock domain assignment resulting from variable assignments or
            evaluations. Once a variable clock domain has been established
            it cannot be changed and all further accesses to that variable
            must adhere to the established clock domain(s).
            
            For selectvalue, static, queue and priority mode variables the value
            to be assigned must either be the same clock domain as the current
            clock domain or must have no clock domain (external inputs with no
            clock domain or constants). The clock domain of the assigned
            variable (write clock in the case of queues) must not differ from
            the current clock domain - if it has no clock domain the
            assignment will set it to the current domain.
            
            Note that where variables, or expressions containing variables, of
            mode selectvalue, static, queue or priority are assigned to value
            variables, the output clock domain of the RHS variables is not
            resolved until the value variable is actually directly used. The
            same consideration applies to arguments passed to module, procedure
            or function parameters. Clock domain resolution of variables is
            delayed as late as possible until they are actually assigned to or
            their values are finally required. Their clock domains remain
            unresolved through assignment to value variables or argument
            passing.
            
            Where an assigned value violates the clock domain it may be converted
            by the inbuilt function {\bf sample()} section \autorefph{sec:ifsample}
            or acceptance may be forced by the inbuilt function {\bf accept()}
            \autorefph{sec:ifaccept}. {\bf accept()} simply makes its argument
            appear to be in the current clock domain (or another clock domain or no
            clock domain) but otherwise has no effect. {\bf sample()} resamples its
            argument into the current clock domain. The use of {\bf accept()} may
            result in data corruption unless it is known that the clocks are
            suitably related in phase and frequency. The arguments to {\bf
            accept()} and {\bf sample()} must not contain queue reads.

            A queue which has identical read and write clock domains is termed
            a {\bf synchronous queue}. If it has different read and write
            clock domains it is called an {\bf asynchronous queue}.

            For a queue variable the clock domain for writing is established by -
            \index{queue!availability}
            \blist
            \item   the first write (assignment) to the queue variable (or
                    component of the variable if the type is compound)
            \item   the use of the write availability operator \verb'>' on the queue
            \item   a call to {\bf queuewritestatus()} on the queue
            \item   a call to {\bf reset()} on the queue
            \blend
            If the write clock domain is not yet established in any of these
            situations then it will be set to the current clock domain,
            otherwise it must be the same as the current clock domain.

            The clock domain for queue reading is established by -
            \blist
            \item   the first read or examine of the variable (or component of
                    the variable if the type is compound)
            \item   the use of the read availability operator \verb'<' on the queue
            \item   a call to {\bf queuereadstatus()} on the queue
            \blend
            If the read clock domain is not yet established in any of these
            situations then it will be set to the current clock domain,
            otherwise it must be the same as the current clock domain.
        
            The inbuilt procedure {\bf domain()} may be used to explicitly set
            the clock domain of a selectvalue, static, queue or priority variable
            or a value variable which is an external input. On calling {\bf
            domain()} the clock domain of the variable must be null (in which
            case it is set) or else must already have that clock domain. In the
            case of a queue variable the input and output clock domains will both
            be set, i.e. it will be a synchronous queue.
            
            The inbuilt procedure {\bf readdomain()} may be used to explicitly
            set the output (read) clock domain of a queue variable (for other
            modes it functions the same way as {\bf domain()}).
            
            The inbuilt procedure {\bf writedomain()} may be used to explicitly
            set the input (write) clock domain of a queue variable.

            Memory variables have a clock associated with each access port.
            The clock associated with a port is determined by the clock that
            is current when the first read or write procedure or function
            call for that port is made during immediate execution. All
            subsequent calls to read or write procedures or functions for
            that port must be in the same clock domain.
                      
            A clock variable may be used effectively as a logical variable,
            but this can only be done by assigning it to a value mode
            variable of type log and then that value variable can be used
            in the normal way. The assignment may be explicit, e.g.
            \begin{lstlisting}[gobble=12, language=threepl]
               clock(clk, ....);
               value(v, "log");
               v := clk;
            \end{lstlisting}
            or it may be done via initialisation
            \begin{lstlisting}[gobble=12, language=threepl]
               clock(clk, ....);
               value(v, "log", clk);
            \end{lstlisting}
            
            A clock may be driven by internal logic by use of the inbuilt
            procedure {\bf clockfromexpr()}, e.g.
            \begin{lstlisting}[gobble=12, language=threepl]
               clock(clk, frequency=....);
               value(v, "log");
               clockfromexpr(clk <- v);
            \end{lstlisting}
            
            If an internal clock signal is to be connected to an output pad
            this can be done by passing the clock variable as a source
            argument (third argument) to inbuilt procedure {\bf output()}
            \autorefph{sec:ipoutput}.

        \subsection{checking clock domains}
        
            The current clock domain and the previous clock domain can be
            checked by inbuilt functions currentclock() \autorefph{sec:ifcurrentclock}
            and prevclock() \autorefph{sec:ifprevclock} respectively. The returned
            clock domain can be checked against a known clock, e.g.
            \begin{lstlisting}[gobble=12, language=threepl]
               if (currentclock() != c100)
                   .......
            \end{lstlisting}
            can be checked for null, e.g.
            \begin{lstlisting}[gobble=12, language=threepl]
               if (currentclock() == null)
                   .......
            \end{lstlisting}
            or can have its identifier printed -
            \begin{lstlisting}[gobble=12, language=threepl]
               println(identifier(currentclock()));
            \end{lstlisting}
            The returned clock can also be assigned to a clock variable -
            \begin{lstlisting}[gobble=12, language=threepl]
               clock(c);
               c := prevclock();
            \end{lstlisting}
            but note that a null assignment is not allowed.
            
            The clock domain of variables can be determined by inbuilt
            functions domain() \autorefph{sec:ifdomain}, readdomain() \autorefph{sec:ifreaddomain}
            and writedomain() \autorefph{sec:ifwritedomain}. These will return null if
            the domain has not yet been established. The domain can be
            tested, printed etc as illustrated in the previous paragraph.

        \subsection{setting up clocks}

            \index{clock!source}
            Clock devices are highly part-specific. In general implementing a
            clock requires the following -
            \blist
            \item   The clock signal and any intermediate signals must be
                    declared to 3PL and connected via an FPGA input buffer
                    using the inbuilt procedure {\bf clock()}.
            \item   A clock driver such as a DLL or DCM must be created. This
                    is done via library procedures (see below).
            \item   An FPGA global clock buffer must be provided, for which
                    there is a library procedure {\bf clockbufg()}.
            \item   Additional timing constraints between clock groups may be
                    needed, entered by using the library procedure {\bf
                    maxdelay()}.
                    \index{timing constraint}
                    \index{constraint!timing}
            \blend

            The inbuilt procedure {\bf clock()} is used to create clock
            variables, including any input and intermediate clock signals
            associated with the clock generation hardware in the FPGA.
            The frequency or period attributes must be supplied
            and optionally the duty cycle attribute can be specified. For an
            existing clock these parameters can be determined by the
            inbuilt functions {\bf frequency()}, {\bf period()} and
            {\bf dutycycle()}.
            
            Library procedures clkdll(), clkdlle(), clkdllhf(), clkdcm(), clkmmcm()
            and pll() create the hardware elements which take an input or
            internal clock and produce a related output clock. That output clock
            may be simply the input clock conditioned by a delay-locked loop to
            achieve zero delay or it may be a multiple, submultiple or
            phase-shifted version of the input clock. The arguments passed to
            this procedure and the capabilities it accesses differ with the FPGA
            family. It can usually be called via library procedures which offer
            a simpler interface for common situations.
        
            By default if the clock variable is used for any target
            statements a timing group identifier will be generated for the clock
            variable. This identifier is used to automatically generate a clock
            period timing constraint for all logic driven by the clock. If the
            clock variable has global scope the timing group identifier will be the
            variable identifier converted to upper case. For any other scope the
            timing group identifier will be the full hierarchical variable name
            converted to upper case. The latter case is necessary to avoid any
            possible ambiguity but results in a long identifier, and this is
            impossible to use in any other timing constraint since it has arbitrary
            integers embedded to disambiguate specific instances. Where a clock
            variable is not global and its timing group identifier is to be used in
            generating other timing constraints, e.g. between clock groups, a
            timing group identifier must be supplied to replace the default id.
            This is done by setting the attribute {\bf tnm} in the
            clock variable.
            
            Note that the intermediate clock signals interconnecting the input
            buffers, clock managers and global clock buffers are all variables
            of mode {\bf clock} and potentially can all be used by 3PL to
            control logic. However, Clock variables that can actually be used to clock
            logic devices are restricted to those driven by
            appropriate FPGA devices such as global clock buffers. The
            limitations on which clock signals can be used to drive logic are
            FPGA specific and 3PL itself makes no distinction.
            
            Where clock variables are assigned using the \verb':=' operator
            and are also used to set the clock domain, note that logic
            using that clock will be connected to the clock source determined
            by the value the clock variable has when immediate execution
            terminates. Changing the value of a clock variable by assignment,
            even if accompanied by passing it to {\bf useclock()}, will not
            result in associated code having different clock domains - the code
            will all be driven by the clock source given by the final value of
            the clock variable.

            Further details may be obtained in the sections on the header files
            for the individual boards.
















    \section{Modules, procedures, functions and classes}
        

        \subsection{Module, procedure, function and class calls}\label{sec:mpfccall}
        
            Procedures, named modules and classes are called by
            \begin{lstlisting}[gobble=12, language=threepl]
               name (output arguments <- input arguments);
            \end{lstlisting}
            and if there are no output parameters -
            \begin{lstlisting}[gobble=12, language=threepl]
               name (input arguments);
            \end{lstlisting}
            
            Functions have no output parameters and are called by
            \begin{lstlisting}[gobble=12, language=threepl]
               name (arguments)
            \end{lstlisting}
            A function call appears within an expression (it may be the entire
            expression itself).
            
            When a class, named module, procedure or function is called
            the current scope, file scope and global scope are searched
            in that order.
            
            Class calls are similar to function calls, the returned result being
            a class instance which is an immediate value of type "class".
            Unlike functions class calls may contain output arguments.
        
            For classes, procedures, named modules, functions or classes if
            there are no arguments the parentheses are present but empty.
            
            Note that the separation of arguments into input and output
            groupings is mostly a convention, but for some parameter
            modes the mechanism by which an argument is passed to a parameter
            differs depending on whether it is in the input group or the output
            group (see \autorefph{sec:callparams} below.

            \index{argument!key match}
            Arguments are matched against parameters by position in the
            argument list, but for declared modules, procedures or
            functions (i.e. not inbuilt ones) they may also be paired
            explicitly by using
            \begin{lstlisting}[gobble=12, language=threepl]
               parameter = argument
            \end{lstlisting}
            for each argument in the list where {\em parameter} is the parameter
            identifier. This is refered to as "key match" argument passing. The
            effect of using the '=' to match an argument to a parameter is simply
            to re-order the argument list so that the arguments appear in the order
            of the parameters, arguments being omitted where parameters have not
            been matched. This means that arguments will be evaluated and assigned
            to parameters in parameter order, left to right, not necessarily in
            argument order.
            
            Key-match argument/parameter matching is useful where arguments are
            sparse since it avoids runs of commas and resulting mistakes in
            positioning arguments. It can also make code more readable by giving
            clues to the purpose of each argument, but argument lists then
            consume more space.
                
            Null arguments cannot appear among key-match arguments. If they do
            the fatal error "too many input arguments" will occur, or else in
            the case of a {\bf varargs} parameter (see \autorefph{sec:varargs})
            "null argument to input varargs" or "null argument to output
            varargs" errors.
            
            The constant {\bf null} may be passed to a parameter. If the
            parameter is immediate mode and its type is "\verb'->'", i.e. it is a
            pointer, then the pointer is assigned the null pointer. For
            any other parameter mode and type the assignment is skipped.
            Thus except for the case of pointers, the argument {\bf null}
            can be used to conditionally skip passing an argument to a
            parameter. This applies both to argument passing by position or
            by key=value.

            Non-key-match arguments cannot appear among key-match arguments unless
            the last parameter is {\bf varargs} (see
            \autorefph{sec:varargs}). If they do the fatal error "too many
            input arguments" will occur.
            
            Argument matching to parameters by position and by key match can be
            mixed in one call but position matched arguments must precede any
            key matched arguments, i.e. no position matched arguments may appear
            after a key matched argument (unless {\bf varargs} is used). The
            occurence of the first key-match argument in a list signals the end of
            matching arguments to parameters by position.
            
            Key-match argument passing cannot be used with inbuilt modules,
            procedures or functions at the moment since these do not have
            parameters as such. This may be changed in the future.
            
            Key-match argument passing  can also be used to pass arguments
            to a {\bf varargs} list - see \autorefph{sec:varargs}.

            Procedure and function calls may be recursive\index{recursion}.
            Because calls are made during immediate execution, not target
            execution, recursion does not involve code in the FPGA. In principle
            class and module calls can also be recursive but it is hard to
            envisage a case where this would be useful.

        
        
        
        
        \subsection{Module, procedure, function and class declarations}\label{sec:mpfcdecl}
        
            \index{module!declaration}
            A named module is declared by -
            \begin{lstlisting}[gobble=12, language=threepl]
               module
               name (output parameters <- input parameters) 
                   body
               }
            \end{lstlisting}
            or
            \begin{lstlisting}[gobble=12, language=threepl]
               module
               name (input parameters) {
                   body
               }
            \end{lstlisting}
            if there are no output parameters. Similarly a procedure
            declaration appears as -
            \begin{lstlisting}[gobble=12, language=threepl]<<=
               procedure
               name (output parameters <- input parameters) {
                   body
               }
            \end{lstlisting}
            or
            \index{procedure!declaration}
            \begin{lstlisting}[gobble=12, language=threepl]
               procedure
               name (input parameters) {
                   body
               }
            \end{lstlisting}
            A function is declared by -
            \begin{lstlisting}[gobble=12, language=threepl]
               function
               name (parameters) {
                   body
               }
            \end{lstlisting}
            and has no output parameters.
            
            A class is declared by -
            \begin{lstlisting}[gobble=12, language=threepl]
               class
               name (output parameters <- input parameters) 
                   body
               }
            \end{lstlisting}
            or
            \begin{lstlisting}[gobble=12, language=threepl]
               class
               name (input parameters) {
                   body
               }
            \end{lstlisting}
            if there are no output parameters.
             
            \index{module!leadermodule}
            \index{leadermodule}
            A leader module is declared by -
            \begin{lstlisting}[gobble=12, language=threepl]
               leadermodule (s, n) {
                   body
               }
            \end{lstlisting}
            Argument {\bf s} is a comment string which is printed in the report
            file when the module is executed. Optional argument {\bf n} is an
            ordering integer. If omitted its default value is 0. The new module
            will be placed in the leader module list after any modules with the
            same or lower order number.
            
            {\bf Note that for leader, trailer or post-processing modules both
            arguments must be literal constants, not variables or expressions, as
            they are extracted from tokens during parsing, not evaluated during
            immediate execution.}

            Since it is not called directly it has no identifier. As each leader
            module declaration is encountered during parsing it is added to the
            leader module list. Modules in this list are executed in order prior
            to beginning normal immediate execution.
            
            \index{module!trailermodule}
            \index{trailermodule}
            A trailer module is declared by -
            \begin{lstlisting}[gobble=12, language=threepl]
               trailermodule (s, n) {
                   body
               }
            \end{lstlisting}
            Argument {\bf s} is a comment string which is printed in the report
            file when the module is executed. Optional argument {\bf n} is an
            ordering integer. If omitted its default value is 0. The new module
            will be placed in the trailer module list after any modules with the
            same or lower order number.

            Since it is not called directly it has no identifier. As each
            trailer module declaration is encountered during parsing it is added
            to the trailer module list. Modules in this list are executed in
            order after normal immediate execution terminates.
            
            If a trailer module is defined within a class then the class
            variables, common block variables, modules, procedures and functions
            will be accessible to the trailer module.
            
            \index{module!postmodule}
            \index{postmodule}
            A post-processing module is declared by -
            \begin{lstlisting}[gobble=12, language=threepl]
               postmodule (s, n) {
                   body
               }
            \end{lstlisting}
            Argument {\bf s} is a comment string which is printed in the report
            file when the module is executed. Optional argument {\bf n} is an
            ordering integer. If omitted its default value is 0. The new module
            will be placed in the post module list after any modules with the
            same or lower order number.

            Since it is not called directly it has no identifier. As each
            post-processing module declaration is encountered during parsing it
            is added to the post-processing module list. Modules in this list
            are executed in order after netlist and constraint file generation
            has been completed. This code can invoke place-and-route
            tools and any other post-processing operations which can again
            be conditional on the time of last modification of related files.
            
            See also \autorefph{sec:lppostprocess} and \autorefph{sec:xpp}.
                     
            Module, procedure, function and class declarations may be nested. A
            module, procedure, function or class can be called from within the
            module, procedure, function or class in which it is declared or from
            any nested module, procedure, function or class.

            In the case of declarations in a file module the keywords {\bf
            module}, {\bf procedure}, {\bf function} or {\bf class} may be
            preceded by the keyword {\bf global} which makes the module, procedure
            or function globally accessible. The keyword {\bf global} cannot be
            used on declarations nested deeper than the file scope.
        
            \index{identifier}
            Identifiers fall into four categories -.
            \blist
            \item   variables
            \item   directives
            \item   modules, procedures and classes
            \item   functions and classes
            \blend
            The same identifier may be used in each of these categories as the
            syntax disambiguates the multiple occurrences. The overlap with
            classes is because a class is called in a similar manner to a
            function but can also be called as for a procedure, the returned
            value being ignored. It is thus preferable identifiers for modules,
            procedures, functions and classes are mutually unique.
            
            The same identifier may also be used in different scopes in which
            case the search order determines which is used. In the case
            of modules, procedures and functions the most deeply nested will
            be found first.

            For procedures, named modules, functions or classes if there are no
            parameters the parentheses are present but empty.
            
            A {\bf leadermodule} or {\bf trailermodule} can create variables in
            the global scope (global.{\em name}) or in the file scope in which
            it is declared (file.{\em name}).

        \subsection{parameters and arguments}\label{sec:callparams}
        
            \index{module!parameter}
            \index{procedure!parameter}
            \index{function!parameter}
            \index{argument!passing}
            Parameters are declared using the same inbuilt variable creation
            procedure calls as normal variables (they cannot be user-defined
            procedures). These procedure calls are enclosed in the parentheses
            following the module, procedure or function identifier. The type of a
            parameter must be the same as that of the corresponding call argument
            but word widths need not be the same. Where word widths differ the
            assignment rules described previously (truncation or zero/sign padding)
            apply.

            Parameter variables which are not matched by corresponding
            arguments become local variables. Parameters can be initialised so
            that they have useful default values should they have no matching
            arguments in a call. While obviously useful for input parameters
            this also applies to output parameters.

            Parameters may be any mode.

            The argument to an immediate input parameter can be
            an immediate expression, which includes the cases of a single
            immediate variable or constant. An immediate mode parameter is always
            assigned an argument by value, even if the argument is itself
            an immediate variable, so if an immediate input parameter
            is assigned within a module, procedure or function its local value
            is simply changed and the argument is not affected.

            The argument to a value input parameter can be an expression,
            including the cases of single target or immediate variables or
            constants. An immediate argument will be converted to a target
            constant. Value input parameters are call-by-value if the argument
            is not itself a single value variable. If the argument is a single
            value variable the argument is passed by reference.
            
            A value input parameter may be passed an immediate compound
            type in which case the parameter must be declared with a matching
            target compound type. Any immediate value which exceeds the
            size of the corresponding target type will cause a fatal error.
            
            A value argument may be passed to an input, selectvalue, static,
            queue or priority output parameter. This causes the argument value
            variable to refer to the parameter variable which itself is now a
            local variable.

            Precluding the cases listed above -
            \blist
            \item   immediate mode input parameter
            \item   value mode input parameter passed anything except a value
                    mode variable
            \item   an output parameter which is input, selectvalue, static,
                    queue or priority mode and is passed a value mode variable
            \blend
            which are call-by-value, all other parameters are
            call-by-reference and can only be passed variables of the same
            mode as arguments. Parameters which are Call-by-reference actually
            use the argument variables instead of the parameter variables.
            
            Where a parameter would normally be assigned an argument value by
            value rather than by reference, as described above, passing by
            reference can be forced by using the inbuilt procedure {\bf ref()}
            \autorefph{sec:ipref} instead of {\bf var()}, {\bf int()}, {\bf
            str()} etc. In that case an argument which is an expression rather
            than a variable will cause a fatal execution error.
            
            Note that the separation of parameters into input and output groups
            is mostly a convention intended to make the code more readable.
            Apart from the differences in the way arguments are passed to
            parameters described above, there are no other semantic differences
            between input and output parameters. In particular input
            parameters may be assigned values (but note that the effect depends
            on whether they are call-by-value or call-by reference) and 
            output parameters may be used in expressions. However, it is preferable
            to adhere to this scheme where possible. An exception is
            where a variable is to be passed to a module or procedure within
            which code appears both to assign to it and to use its value. In that
            case it is best to pass it to an output parameter.

            A value parameter is a very general way of passing target
            arguments to a module, procedure or function. The argument to a
            value parameter can be a constant (or compound value), an
            expression or a variable of any mode except {\em memory} (an
            immediate variable or expression will pass a constant). In the case
            of a floating point target mode argument to a value parameter the
            parameter type must either match the argument type exactly or else
            must have the type omitted. Target floating point types cannot have
            the mantissa or biased exponent sizes omitted and hence no
            adjustment to the argument type is possible. If the parameter type
            is given it must match the argument type exactly. If the parameter
            type is missing it will take the type from the argument. In the
            case where the argument is an immediate floating point type it will
            be converted to a target type "float:23.8" when passed in the call.

            An argument to a call-by-reference parameter may be a partial
            variable reference, for example an array member, a sub-array or a
            struct field. This is allowed as long as there is no type mismatch
            with the associated parameter. If the parameter is not typed or has
            an array dimensions omitted then the type or array dimensions can
            be taken from the argument. The use of a partial variable as an
            argument applies to the queue mode as well as all other modes.

            Wherever a call-by-reference parameter is used in the code body,
            the argument variable that it points to will be used instead. If a
            parameter has not been passed an argument then it becomes simply
            a local variable within the scope of the module, procedure or
            function.
            
            Since call-by-reference parameters use the argument variable in
            place of the parameter variable the clock domain of the
            argument variable is relevant. Similarly a value parameter inherits
            the clock domain or domains of its argument (a compound
            value may have elements of differing clock domains).
            
            Note that it is important in which module a queue is read. This is
            regardless of whether it is passed as an argument to a module,
            procedure or function where it appears as a parameter. Ultimately a
            queue variable, even if referenced via a parameter (or chain of
            arguments/parameters), is read within a module.
            Each module in which a queue is read constitutes a queue destination.
            Each queue destination {\bf must} read a datum before it can be
            popped from the queue. When a destination module reads a datum from a
            queue the queue will appear to it to be unavailable unless and until
            all other destination modules execute a read (they can of course all
            read simultaneously). Reads in multiple locations in a destination
            module read the same datum if they execute in the same clock cycle,
            otherwise they will read different data.            

            Where a parameter is static or queue mode the parameter will
            have the clock of the matching argument variable.


            \subsubsection{missing parameter target primitive widths}
            
                \index{parameter!missing target width}
                A target mode parameter which is type {\em bits}, {\em uint} or
                {\em int}, or is an array of one of these types, may have the
                width missing following the colon, or may have width 0. In that
                case the type match with the supplied argument will be checked
                and the width will be taken from the argument. If the parameter
                variable width is supplied it must agree with the width of the
                argument variable. Note that the colon is still necessary in
                the type string to indicate a target mode type. If the argument
                is a constant the width will be the minimum necessary to
                contain the constant taking into account the sign bit if type
                {\em int}.\footnote{At the moment a target parameter array
                variable cannot derive the type width from an argument which is
                a compound value. In that case the target variable width
                will be zero, giving a fatal error. This may be fixed in the
                future.}        

            \subsubsection{missing parameter dimensions}

                \index{parameter!missing dimension}
                A parameter type specification may have missing array
                dimensions (empty '[]') in which case the type match with the
                supplied argument will be checked and missing dimensions will
                be copied from the argument type. Note that this cannot be
                done if the parameter is created as one of a bracketed list,
                e.g.
                \begin{lstlisting}[gobble=16, language=threepl]
                   procedure
                   p (var({a, b}, "[]int")) {
                   }
                \end{lstlisting}
                is not allowed as both {\bf a} and {\bf b} must be the same type,
                but the array dimension is not supplied and the arguments may be of
                different dimensions. What is allowed is a common initialiser, e.g.
                \begin{lstlisting}[gobble=16, language=threepl]
                   procedure
                   p (var({a, b}, "[][]int", {{3, 5, 1}, {4}})) {
                   }
                \end{lstlisting}
                as here both {\bf a} and {\bf b} will be two-dimensional arrays
                of type "[2][3]int" and both will be initialised.
                
                Where a parameter type has one or more missing dimensions and
                both an initialiser and an argument are supplied the dimension
                of the argument will be used. If an argument is not supplied
                the dimension of the initialiser will be used.

            \subsubsection{missing parameter types}

                \index{parameter!missing type}
                If a parameter is declared without a type at all, the type of
                the call argument will be used.
                A parameter variable with missing type must be created singly,
                not as one of a bracketed list.

            \subsubsection{missing parameter modes}

                \index{parameter!missing mode}
                If a parameter is declared without a mode, the mode of
                the call argument will be used.
                A parameter variable with missing mode must be created singly,
                not as one of a bracketed list.

            \subsubsection{type conversions}
            
                \index{string/type conversion}
                In the particular case of a parameter of type "str" and argument of
                type "type" the argument will be converted to string representation
                of the type (the same as invoking unbuilt function {\bf
                typestring()}). For a parameter of type "type" and argument of type
                "type"  the string will be converted to a type.
                
                The same conversion occurs where a variable is created as one of
                these types and an initialisation argument is of the other type.

            \subsubsection{parameter/argument processing order}

                \index{argument!passing!order}
                For procedures and named modules the input parameter list is
                processed first followed by the output parameter list.
                Parameter lists are processed left to right. Within this order
                later parameters can use characteristics derived from those
                already processed.

            \subsubsection{parameter testing}

                \index{parameter!checking}
                Parameters can be checked to determine if they were supplied an
                argument by the use of the inbuilt function {\bf matched()}.
                Undimensioned array parameters can have their dimensions
                determined after argument matching by using the inbuilt
                functions {\bf dimensions()} and {\bf dimension()}. Untyped
                parameters can have their types determined after argument
                matching by using the inbuilt function {\bf type()};

            \subsubsection{variable number of arguments}\label{sec:varargs}

                \index{argument!variable number}
                \index{varargs}
                Where the input or output arguments are unknown in number or mode
                the input or output parameter list (or both) may contain the
                keyword {\bf varargs} to match a number of arguments. {\bf
                varargs} must be the last item in an input or output parameter
                list i.e. it cannot be followed by any variable
                creation procedure calls.

                When {\bf varargs} appears in a parameter list
                local array variables {\bf inargs} and {\bf outargs} as
                appropriate will be created automatically. For each argument an
                invisible parameter of appropriate mode will be created. For
                arguments which are variables the mode will be the same as the
                variable. For a target expression the parameter variable will
                be mode {\em value}. These {\bf inargs} and {\bf outargs} arrays are
                pointers to the created invisible parameters. Immediate code
                can interrogate these invisible parameters via the pointer
                arrays to determine mode, type, dimensions etc. using inbuilt
                functions {\bf mode()}, {\bf type()}, {\bf typestring()}
                (and possibly {\bf dimensions()} and {\bf dimension()}, though
                {\bf typestring()} returns this information in string form).
                Inbuilt function {\bf dimension()} can be used to get the
                number of arguments from {\bf inargs} or {\bf outargs}. The
                invisible parameters are accessed via the pointers, e.g.
                \begin{lstlisting}[gobble=16, language=threepl]
                    a <- inargs[1]->;
                    outargs[3]-> <<- a + b;
                \end{lstlisting}
                If an argument has fields or array members then field
                identifiers or array subscripts can follow the '\verb'->'',
                e.g.
                \begin{lstlisting}[gobble=16, language=threepl]
                    a <- inargs[1]->[2];
                    outargs[3]->data <<- a + b;
                \end{lstlisting}
                
                \index{varargs!inargs}
                \index{varargs!outargs}
                \index{varargs!inkeys}
                \index{varargs!outkeys}
                \index{varargs!inmap}
                \index{varargs!outmap}
                A function, module or procedure whose parameter list
                contains {\bf varargs} can be passed key-match arguments.
                Where the key-match arguments identify explicit parameters
                the arguments will be assigned to the associated parameter
                variables as usual. If the key does not match a parameter
                identifier the key-argument pair will be assigned to the
                varargs list. For this reason the varargs pointer arrays
                {\bf inargs} and {\bf outargs} have associated arrays
                {\bf inkeys} and {\bf outkeys} which are arrays of strings,
                and {\bf inmap} and {\bf outmap} which are maps containing
                pointers to arguments addressed by the key string.
                
                String arrays {\bf inkeys} and {\bf outkeys} contain the keys
                from key-matched arguments. Where an argument to varargs is not
                key-matched the associated key array entry is "". These arrays
                of keys are always present when {\bf varargs} is used but if
                there are no key-matched arguments then the key arrays will
                contain empty strings.
                
                In addition to the above, key-match arguments are also entered
                into maps, {\bf inmap} and {\bf outmap}, where the argument key
                is the map key and the associated value is a pointer to the
                argument value.
                
                Maps {\bf inmap} and {\bf outmap} only contain key-match
                entries, not positional arguments. Map access using a key that
                is not present in the map will return a null.                
                
                As an example, the procedure -
                \begin{lstlisting}[gobble=16, language=threepl]
                   procedure
                   p (int({a, b}), varargs) {
                   }
                \end{lstlisting}
                might be called by -
                \begin{lstlisting}[gobble=16, language=threepl]
                    p(, 2, k=3, a=5, n=4);
                \end{lstlisting}
                The first parameter {\bf a} is skipped and the second parameter
                {\bf b} is positionally assigned 2. The next three arguments are
                key-match, but the second of these will actually assign 5 to
                parameter {\bf a}. Thus the varargs array {\bf inargs} will contain
                pointers to immediate integer variables containing 3 and 4 and the
                corresponding array {\bf inkeys} will contain "k" and "n".
                
                Note that we could also have non-key-match arguments (but not null
                arguments) among the key-match arguments passed to varargs, e.g.
                \begin{lstlisting}[gobble=16, language=threepl]
                    p(2, k=3);

                    procedure
                    p (varargs) {
                    }
                \end{lstlisting}
                would result in array {\bf inargs} containing pointers to
                immediate integer variables containing 2 and 3 and the
                string array {\bf inkeys} will contain "" and "k".
                
                If we did not have {\bf varargs} as the last parameter then
                non-key-match arguments and null arguments cannot appear
                among key-match arguments. If they do the fatal error
                "too many input arguments" will occur.
                
                A compound value cannot be passed to a {\bf varargs}
                parameter as there is no type available with which to
                parse the compound value.
                
                These argument passing mechanisms are intended to provide
                \blist
                \item   argument matching to parameter by position
                        or parameter identifier
                \item   variable length, type and mode argument list
                \item   extension of key=argument to variable argument lists
                \blend
                            
                It is envisaged that module, procedure and function calls
                would use one of the following scenarios for fixed parameter
                lists, i.e. {\bf varargs} not used -
                \blist
                \item   positional arguments only
                \item   key-match arguments only
                \item   first few arguments positional, the remainder key-match
                \blend

                For variable argument lists, i.e. where a {\bf varargs} parameter
                is used -
                \blist
                \item   parameters with final varargs - positional matching
                        arguments to parameters, non-key-match and key-match
                        arguments to varargs
                \item   varargs parameter only - non-key-match and key-match
                        arguments
                \blend
                The situation where some key-match arguments match parameters
                and others are assigned to varargs and some positional
                arguments match parameters and yet others are assigned to
                varargs is quite legal but is probably much too confusing.
                
                Variable arguments can be passed on to nested calls by using
                inbuilt function {\bf arglist()} \autorefph{sec:ifarglist}.
                Note however that this cannot be used for calls to inbuilt modules,
                procedures or functions.
                
                As an example -
                \begin{lstlisting}[gobble=16, language=threepl]
                    procedure q1 (varargs) {
                        q2(arglist(null, inargs));
                    }
                \end{lstlisting}
                In this case the input arguments to procedure q1 will be passed
                to procedure q2. If inargs is simply passed directly it will
                appear as a one-dimensional array of pointers. By passing
                via arglist() the values pointed to will appear as separate arguments.
                

        \subsection{target code}

                \index{immediate!target generation}
                During immediate execution modules, procedures, functions and
                classes exist as single code bodies even though perhaps called
                repeatedly. In that sense they function in the same way as normal
                programming languages. However, any target code that they generate
                is not shared in any way. Target code is instantiated afresh on
                every call.

                For example
                \begin{lstlisting}[gobble=16, language=threepl]
                    static(a, "[4]int:8");
                    static(b, "[4]int:8");
                    var(double, "[4]int");

                    for (int(i) ; i<4 ; i++)
                        reverse(a, double[i] <- b, i);

                    procedure
                    reverse (static(o, "[]int:8"), int(d) <- static(i, "[]int:8"), int(j) {
                        d = j * 2;
                        o[j] <- i[3-j];
                    }
                \end{lstlisting}
                Here an immediate loop calls a procedure four times. Each time
                the immediate integer argument {\em i} is doubled and stored in
                the argument array {\em double}. However, the procedure also
                generates a target assignment on each call. It is important to
                note here that we have generated four target assignment
                statements as a result of immediate execution. It is as though
                we wrote -
                \begin{lstlisting}[gobble=16, language=threepl]
                    double[0] = 0 * 2;
                    a[0] <- b[3];
                    double[1] = 1 * 2;
                    a[1] <- b[2];
                    double[2] = 2 * 2;
                    a[2] <- b[1];
                    double[3] = 3 * 2;
                    a[3] <- b[0];
                \end{lstlisting}
                Execution control of these target statements must be considered.
                As written the above code must be immediately within a module to
                be correct. Each of the generated target assignment statements
                will be wrapped in its own infinite loop to execute on every
                clock cycle (see next section). However, suppose that the
                immediate loop above is contained within a {\em when} statement.
                In that case we must enclose the four target assignments in a
                parallel or sequential block, e.g.
                \begin{lstlisting}[gobble=16, language=threepl]
                    static(a, "[4]int:8");
                    static(b, "[4]int:8");
                    var(double, "[4]int");

                    when (l) par {
                        for (int(i) ; i<4 ; i++)
                            reverse(a, double[i] <- b, i);
                    }

                    procedure
                    reverse (static(o, "[]int:8"), int(d) <- static(i, "[]int:8"), int(j) {
                        d = i * 2;
                        o[j] <- i[3-j];
                    }
                \end{lstlisting}
                which is equivalent to -
                \begin{lstlisting}[gobble=16, language=threepl]
                    when (l) par {
                        double[0] = 0 * 2;
                        a[0] <- b[3];
                        double[1] = 1 * 2;
                        a[1] <- b[2];
                        double[2] = 2 * 2;
                        a[2] <- b[1];
                        double[3] = 3 * 2;
                        a[3] <- b[0];
                    }
                \end{lstlisting}
                Here immediate execution of statements within the {\em par} is
                carried out sequentially. Each immediate assignment is carried out
                as it is encountered. Each target assignment is appended to the
                parallel block code body to be eventually executed in parallel in
                the FPGA. The {\bf when} itself will be within an infinite loop.






    \section{Procedures}

        \index{procedure}
        A procedure is simply a way of writing code which can be used in one
        or more places with the calling arguments effectively substituted for
        the parameters. It also provides a way of hiding
        variables which would otherwise be duplicated (variables local to the
        procedure).

        During immediate execution a procedure may execute any number of immediate
        or target statements, but it must only create at most a single in-line
        target statement. If a procedure created more than one in-line target
        statement the behaviour of these statements would depend on the calling
        context, i.e. the behaviour would be different if called from the top level
        inside a module (each statement would then be in its own infinite loop) or
        elsewhere (it must be inside a {\bf seq}, {\bf par} or {\bf sync} in which
        case these environments would control the way the procedure body executed.
        To achieve consistency in the way a procedure body behaves the single
        in-line target statement rule has been imposed. The procedure body can now
        explicitly use a {\bf seq}, {\bf par} or {\bf sync} block for the code body
        or it can use an inline module so that the enclosed statements are in
        individual infinite loops.

        Any queues which are used within a procedure, whether they be global or
        file scope or passed via arguments, have the calling module as their
        destination module for signaling purposes.
        
        A procedure may return from within its code block by the
        statement {\bf return;}. If there is no return statement then the procedure
        will return after completing execution.
        








    \section{Functions}


        \index{function}
        A function code block may only contain
        \blist
        \item   immediate statements
        \item   value, selectvalue, pointer or priority mode assignments
        \item   module calls
        \item   inline modules
        \blend
        It may contain procedure or function calls if they also observe the
        same constraints. The reason for this restriction is that a function
        can only contain code which references variables or expressions. It
        cannot contain code which is one or more elements in a target
        execution stream. It is assumed that the function will execute when
        required (subject to any contained queues) and it cannot contribute
        any extra execution cycles. However, the function may  contain one or
        more modules since they each contain one or more independent
        execution streams. The function in that case may simply
        return a value which is a variable, or an expression containing one or
        more variables, modified by the module.
        
        \index{value!function return}
        \index{statement!return}
        A function can return a value of any mode.
        
        A function returns a value by means of the statement {\bf
        return(\dots);}

        For example a function to double a value mode integer could be written -
        \begin{lstlisting}[gobble=8, language=threepl]
            function
            double (value(p, "int:8")) {
                return(p + p);
            }
        \end{lstlisting}
        This simply substitutes the expression for the call with the parameter
        {\em p} replaced by the value of the argument.
        
        A function code body cannot contain in-line executable target code
        statements. The above example could be legitimately rewritten as -
        
        \begin{lstlisting}[gobble=8, language=threepl]
            function
            double (value(p, "int:8")) {
                value(v, "int:8");
                v := p + p;
                return(v);
            }
        \end{lstlisting}
        
        since the use of value mode variables does not generate in-line target
        code. Value mode variables simply represent target mode expressions.If
        however we attempted to change the example to -
        
        \begin{lstlisting}[gobble=8, language=threepl]
            function
            f (value(p, "int:8")) {
                static(a, "int:8");
                a <- p + p; // ERROR! in-line target code
                return(a);
            }
        \end{lstlisting}
        
        This would give a fatal immediate execution error. This contains an
        executable assignment statement which generates in-line target code.
        It is not even clear exactly what this function should do. The
        implication is perhaps that the evaluation of the function, when it
        is called from an expression, will require a delay of one clock cycle
        while the static assignment takes place. Target code execution
        signaling does not allow this to be done this way! Note that removal
        of the static assignment into a procedure would not make any
        difference, e.g.
        \begin{lstlisting}[gobble=8, language=threepl]
            function
            f (value(p, "int:8")) {
                static(a, "int:8");
                x(a <- p);
                return(a);
            }

            procedure
            x (static(o, "int:8") <- value(i, "int:8")) {
                o <- i + i;
            }
        \end{lstlisting}
        
        is still erroneous for the same reason - the procedure call
        effectively just substitutes its body in place of the call so we
        still get the previous code.

        A target mode function can contain in-line target statements if
        it isolates them within a module. That way those statements are
        quite separate from the function body.
        For example, suppose that we wish to use a static variable which has been
        delayed by one clock cycle. We could write a function for that as
        follows -
        
        \begin{lstlisting}[gobble=8, language=threepl]
            function
            delay (value(p, "int:8")) {
               static(a, "int:8");
               [
                   a <- p;
               ]
               return(a);
            }
        \end{lstlisting}
        The function body is now legal as it does not itself contain any in-line
        target statements but generates a separate module which does contain
        such statements. The function simply returns the value of a static
        variable. That static variable is assigned in the independent module
        (a named module call could have been used instead of an un-named module).
        
        Note that in the above example the assignment of the static variable {\bf
        a} is unconditional and thus occurs every clock cycle regardless of when
        the return value of the function is used. In some cases we may need to
        ensure that the assignment only occurs when we use the value of the
        function. This can be done by using the inbuilt function {\bf used()}
        \autorefph{sec:ifused}. This function can only appear in the body
        of a function definition so to use it in an associated module it must
        be passed as a value mode variable to the module -
        
        \begin{lstlisting}[gobble=8, language=threepl]
            function
            f (value(p)) {
                static(a, type(p));
                value(used, "log", used());

                [
                    when (used)
                        a <- p;
                ]

                return(a);
            }
        \end{lstlisting}
        
        As an aside there are two simplifications in this version. Firstly
        the function parameter type has been omitted. It will take the type of
        the argument, which is fine in this case. Secondly the type of static
        mode variable {\bf a} has been copied from the function parameter rather
        than being explicitly given. These make the code more succinct and also
        provide some generality in the function. The type of the argument to
        the function could be any target type however complex.
        
        Note that any source statement containing a call to this function, and
        there may be any number, will create a separate instance of the entire
        function code and its nested module. This is true even if there are
        multiple calls within the same statement; each generates a separate
        instance.
        
        As a note of caution suppose we were to call function {\bf f()} with an
        argument which is a queue or an expression containing queue reads?
        
        \begin{lstlisting}[gobble=8, language=threepl]
            static(s, "uint:8");
            queue(q, "uint:8");
            
            .
            .
            .
            s <- f(q);
            .
            .
        \end{lstlisting}        
        
        The behaviour of this code would probably not be what we would want. The
        module would repeat execution of the {\bf when} until the function value
        is used at which time it will attempt to do the assignment to {\bf a}. If
        all queues in {\bf p} are available (not empty) then the assignment will
        occur and the function will return the value of {\bf a} as before. If
        however all queues in {\bf p} are not available that assignment will wait
        until they are available. Meanwhile the function has no dependence on {\bf
        p} and will simply return the value of {\bf a}. The value of {\bf a} will
        belatedly change when {\bf p} becomes available. An attempt to fix the
        problem using a {\bf sync} will also not work as the queue read is
        actually in a different module to that in which the {\bf sync} appears -
        
        \begin{lstlisting}[gobble=8, language=threepl]
            .
            .
            sync {          // WONT WORK!
                s <- f(q);  // q is read in different module
            }
            .
            .
        \end{lstlisting}
        
        The solution is to abandon the attempt to use a function and a separate
        module and instead simply code it in the following straightforward way -        
        
        \begin{lstlisting}[gobble=8, language=threepl]
            static(s, "uint:8");
            queue(q, "uint:8");
            static(a, type(q));
            
            .
            .
            sync {
                a <- q;
                s <- a;
            }
            .
            .
        \end{lstlisting}
        
        The declaration of a could be moved into the sync block since
        it is not used outside that block -
       
        \begin{lstlisting}[gobble=8, language=threepl]
            static(s, "uint:8");
            queue(q, "uint:8");
            
            .
            .
            sync {
                static(a, type(q));
                a <- q;
                s <- a;
            }
            .
            .
        \end{lstlisting}
        
        and the whole sync block could be placed in a procedure -
        
        \begin{lstlisting}[gobble=8, language=threepl]
            static(s, "uint:8");
            queue(q, "uint:8");
            
            .
            .
            p(s <- q);
            .
            .
            procedure
            p (static(ss) <- queue(qq)) {
                sync {
                    static(a, type(qq));
                    a <- qq;
                    ss <- a;
                }
            }
        \end{lstlisting}

        
        A function call may be followed by subscripts and field identifiers
        where the function returns an arrays, struct or class (arbitrarily nested).
        
        If a function returns a pointer then it may be dereferenced using
        the \verb'->' operator, e.g.
        \begin{lstlisting}[gobble=8, language=threepl]
           a = f()->;
        \end{lstlisting}
        The \verb'->' operator may be followed by subscripts or field identifiers,
        e.g.
        \begin{lstlisting}[gobble=8, language=threepl]
           a = f()->[i].f2;
        \end{lstlisting}
  
        These function call constructions may also appear on the left hand side
        of an assignment statement, e.g.
        \begin{lstlisting}[gobble=8, language=threepl]
           f()->[i].f2 = 5;
        \end{lstlisting}
        
        A function call can also appear as an incrementing or decrementing
        statement. In that case it will be followed by the \verb'->' operator
        and a ++ or --, e.g.
        \begin{lstlisting}[gobble=8, language=threepl]
           f()->[i].f2++;
        \end{lstlisting}
        If the variable pointed to by the function return is static then
        this will generate a target statement.
        
        The example could be further complicated by a multiply indirect
        function call and the return of an array of pointers, e.g.
        \begin{lstlisting}[gobble=8, language=threepl]
           a[2]->b.f1[j]->().f3[4]->[i].f2++;
        \end{lstlisting}        
        







    \section{Modules}\label{sec:modules}

        \index{module}
        A module is a subroutine whose main purpose is to create execution
        threads in the FPGA.
        
        Modules are similar to procedures but differ in two respects -
        \blist
        \item   modules are single destinations for queues
                whereas procedures have no particular significance when
                reading and writing queues
        \item   module top level in-line target statements are each enclosed
                in individual infinite loops, but a procedure can only
                generate one in-line target statement which is inserted
                in-line back at the calling location
        \blend
        Thus modules form the base of target code execution by initiating
        execution threads whereas procedures are simply a means of re-using
        immediate code or improving readability by subdivision. An execution
        thread thus initiated can only terminate if it encounters a {\bf stop()}
        procedure call (\autorefph{sec:ipstop}) around which there is no parallel
        execution path.
        
        Module definitions can be nested but scope rules determine where they can
        be called.

        Module and procedure calls can be nested but executable target statements
        directly within each module have independent infinite loops, i.e.
        execution threads. Where modules are nested each individual module is
        a separate queue destination regardless of the nesting. For example -
        
        \begin{lstlisting}[gobble=8, language=threepl]
            m1(q);
            
            module m1 (queue(q1)) {
                ... <- q1;
                .
                .
                m2(q1);
                .
            }
            
            module m2 (queue(q2)) {
            .
                ... <- q2;
                .
            }
        \end{lstlisting}
        Module m1 is called and in turn calls module m2. Each of the modules
        reads the queue and is thus a destination for the queue. Thus to pop
        an entry from the queue both reads must execute, whether simultaneously
        or at different times.

        Modules usually communicate with each other mainly by means of queues,
        but variables of other target modes can also be used - a module
        need not use queues at all. A module need not contain any target
        executable statements at all, in which case it is effectively a procedure.

        Queues to be read or written by a module may be passed in the call via
        parameters (named module) or else may be in the global scope or in the
        same file scope (file or inline modules). There may be several places in
        a module where a queue is written. These queue writes will either be
        mutually exclusive events or will simultaneously write to different
        parts of a compound type. Writes to words within a queue constitute
        separate queue writes if they occur at different times (clock cycles)
        but are a single queue write if they occur simultaneously (the same
        clock cycle)\footnote{Simultaneous writing to the same word in a queue
        is an undetected execution error. Where this is not precluded by the
        logic it can be avoided by arbitration using the waitpri control
        structure \autorefph{sec:waitpri}}.


        
        Every un-nested target statement within a module is placed within an
        infinite loop. Thus if a file module (source file) contained only -
        \begin{lstlisting}[gobble=12, language=threepl]
            static(s, "uint:8");
           
            s++;
        \end{lstlisting}
        then the only target executable statement, {\bf s++}, would execute on
        every clock cycle. If the statement were nested, say in a {\bf when}
        statement,
        \begin{lstlisting}[gobble=12, language=threepl]
            value(l, "log");
            static(s, "uint:8");
           
            when (l)
            s++;
        \end{lstlisting}
        then the top level statement, the {\bf when}, would execute on every clock
        cycle but the nested increment statement would only execute on the clock
        cycles where {\bf l} is true. Similarly a sequential thread would be created
        by
        \begin{lstlisting}[gobble=12, language=threepl]           
            seq {
                .
                ... target statements
                .
            }
        \end{lstlisting}
        where the sequential block would repeat indefinitely but the enclosed
        target statements would execute sequentially. The sequential execution
        could be initiated by a boolean -
        \begin{lstlisting}[gobble=12, language=threepl]           
            value(l, "log");
            
            when (l) seq {
                .
                ... target statements
                .
            }
        \end{lstlisting}
        While {\bf l} is false the {\bf when} will execute on every clock cycle.
        When {\bf l} is true the {\bf when} will execute the sequential block,
        which will take a number of clock cycles. On completion of the sequential
        block the when will again execute as before.
        
        Any nested modules within a module operate as separate modules,
        independent except for any shared variables.
        
        Another simple type of module is an in-line module. This is simply
        a code block enclosed in brackets, e.g.
        \begin{lstlisting}[gobble=12, language=threepl]           
            [
                .
                ... target statements
                .
            ]
        \end{lstlisting}
        Here again each un-nested statement will execute in its own infinite loop.
        
        An in-line module is "called" when immediate execution reaches it. Thus
        it can be called multiple times (it can even be enclosed in a loop)
        and each occurence produces a separate instance of the target code
        within the module.
        
        A named module is a declared module. It has an identifier and optional
        input and output parameters, e.g.
        \begin{lstlisting}[gobble=12, language=threepl]           
            module m1 (value(out, "uint:16" <- int(i), value(in, "uint:8")) {
                .
                ... statements
                .
            }
        \end{lstlisting}
        Again un-nested statements are enclosed in infinite loops.
        
    

        
        There are six types of module -
        \blist
        \item   file modules are created by a {\bf module} statement which imports
                the source code for the module from a source file. The main
                source file also forms a file module.
        \item   named modules are declared in the source code in a similar way
                to procedures and functions. They are invoked by name.
        \item   inline modules occur as source code enclosed in brackets and are
                called implicitly at the point at which they occur.
        \item   leader modules are created by a {\bf leadermodule()} statement
                and are called automatically once prior to execution of the main
                source code.
        \item   trailer modules are created by a {\bf trailermodule()} statement
                and are called automatically once following completion of
                execution of the main source code.
        \item   post processing modules are created by a {\bf postmodule()}
                statement and are called automatically once following completion
                of execution of any trailer modules.
        \blend
        
        Named modules can be invoked as often as required since they are called
        by name. Inline modules can be invoked more than once
        by being enclosed in a loop. File modules can only be invoked once.

        \subsection{file modules}

            \index{module!file}
            When the main source file is read by 3PL that source file becomes
            a module, called a file module.
            Additional file modules are imported by a module statement, e.g.
            \begin{lstlisting}[gobble=12, language=threepl]
                module "spi.3pl"
            \end{lstlisting}
            A file module can have only one instance, i.e. it can only be imported
            once.
            
            A file module has no parameters and thus cannot be
            passed arguments.
            
            A feature of a file module is that it has its own scope and can thus
            be used to hide variables, limiting their access to code within the
            file module, including nested modules, procedures, functions and
            classes.


        \subsection{named modules}

            \index{module!named}
            Named modules are declared and invoked as described in the
            previous section. Each call creates a new instantiation of a
            module.
            
            A module call will normally be made at the top level with respect to
            target code but this need not be the case. In particular there
            is occasional need to include a module call within a function
            body. For example a function to delay a value by one clock cycle
            cannot have target executable statements in its body; it
            can only return expressions. However, a module to create the delay can
            be called within the function body since the module
            code executes independently. The function would return an
            output value from the module.
            
            It is allowed, but not good practice, to call modules from within
            target structures such as {\bf when}, {\bf branch}, {\bf seq while},
            {\bf par while}, {\bf sync while}, {\bf seq do while}, {\bf par do
            while} or {\bf sync do while}. Since the module target code does not
            execute in-place but in a completely separate execution environment a
            module call within these structures will appear misleading. In
            particular a module call within a target loop structure will not be
            subject in any way to the loop iterations. Note that a module call
            within an immediate loop will indeed be called subject to loop
            iterations.
                        
            A module may return from within its code block by the statement {\bf
            return;} (this includes file modules and inline modules). If there
            is no return statement a module code body simply terminates as usual
            at the end. For the initial file module, whether it executes a
            return or simply executes to the last statement, any trailer modules
            are executed after it terminates.

        \subsection{inline modules}

            \index{module!in-line}
            In-line modules are intended for very simple modules where the
            source code would be more verbose were they to be
            declared and called in the usual way.

            An inline module is written in-line within square brackets,
            i.e.
            \begin{lstlisting}[gobble=12, language=threepl]
                [
                    body
                ]
            \end{lstlisting}
            
            This construct both defines the module and also represents the call,
            i.e. when immediate code execution reaches the in-line module the
            module is 'called'. An alternative way to look at it is that immediate
            code flows into and through the in-line module and if a return statement is
            encountered this is similar to a continue statement in a loop, i.e. it
            skips to the end of the block. The square brackets delineate the code
            block as a separate source or destination for any queue reads or writes
            that it contains, as for all modules.
            
            An in-line module has no parameters and thus cannot be passed
            arguments.

            An inline module is a statement in a code body, but while it
            is allowed to appear within target control structures such as
            {\em when}, {\em sync} etc. it does not make sense to do so
            as the code for the module is not generated in-line but as an
            independent entity. It is normal to include inline modules
            within immediate control structure such as {\em if}, {\em
            for} etc. to make the calls conditional or to replicate them,
            during immediate execution, for example -
            \begin{lstlisting}[gobble=12, language=threepl]
                for (int(i) ; i<6 ; i++)
                    [
                        body
                    ]
            \end{lstlisting}
            which would instantiate six copies of the module. Clearly to make
            sense the content of the module would differ in some way on each
            iteration.

            An inline module may read and write queues. These queues may be created
            outside the module but in an accessible scope or may be created inside
            the module. there are no module parameters or arguments so queues cannot
            be passed to the module by that means.
 
            An inline module has a block scope similar to braced code blocks in
            control structures. The module can access variables back to the
            enclosing procedure, function or named module. It can also access
            its file scope and global scope.
        
            An inline module code block may return conditionally prior to the end
            of the block by the statement {\bf return;}, otherwise it will return
            after the last statement.



        \subsection{leader modules}

            \index{module!leader}
            Leader modules are declared by a {\bf leadermodule()} statement.
            They are created during parsing, not during immediate execution.
            They are invoked prior to the start of immediate execution. An
            optional argument to the statement can be used to determine the
            order of execution of leader modules.
            
            A leader module has two parameters, a comment string and an optional
            order integer. The comment string is used to report execution of the
            leader module in the report file if present. The order integer is
            used to order the execution of multiple leader modules.

        \subsection{trailer modules}

            \index{module!trailer}
            Trailer modules are declared by a {\bf trailermodule()} statement.
            They are created during parsing, not during immediate execution.
            They are invoked following completion of immediate execution of
            the main program. An optional arguments to the statement can be
            used to determine the order of execution of trailer modules.           
            
            A trailer module has two parameters, a comment string and an
            optional order integer. The comment string is used to report
            execution of the trailer module in the report file if present. The
            order integer is used to order the execution of multiple trailer
            modules.
            
            Trailer modules are used to finalise target code that by its
            nature remains incomplete during execution. This is usually
            external interfaces to other devices such as CPUs, memory or other
            FPGAs. For example a multiplexed interface may have a
            number of ports created during immediate execution which are
            unknown in number until main program execution has completed.
            Trailer code can then generate the necessary connecting logic.

        \subsection{post-processing modules}

            \index{module!post-processing}
            Post-processing modules are declared by a {\bf postmodule()}
            statement. They are created during parsing, not during immediate
            execution. They are invoked following completion of immediate
            execution of the main program and any trailer modules. An optional
            arguments to the statement can be used to determine the order of
            execution of trailer modules.
            
            A post-processing module has two parameters, a comment string and an
            optional order integer. The comment string is used to report
            execution of the post-processing module in the report file if
            present. The order integer is used to order the execution of
            multiple post-processing modules.
            
            Post-processing modules are intended to invoke place-and-route
            tools following successful generation of netlist and constraint
            file(s). This can be done by using calls to inbuilt procedures
            {\bf exec()} or {\bf shell()}.
            
            




     \section{Classes}\label{sec:classes}
     
        The class structure is provided for complicated interfaces where multiple
        ports may be added to an interface during immediate execution and the
        logic to tie it all together is added at the end. For most straightforward
        applications classes are not needed and indeed 3PL existed for two decades
        before the need for classes was finally recognised.
       
        A class provides a wrapper around a collection of variables and associated
        modules, procedures, functions and nested classes. In its simplest form
        it provides functionality similar to a struct -
        
        \begin{lstlisting}[gobble=8, language=threepl]
            class c () {
                int(i, 34);
                static(s, "uint:8");
            }
        \end{lstlisting}
        
        An instance of this class can be created by -
        
        \begin{lstlisting}[gobble=8, language=threepl]
            var(cl, "class");
            cl = c();
        \end{lstlisting}
        
        or more simply -
        
        \begin{lstlisting}[gobble=8, language=threepl]
            var(cl, "class", c());
        \end{lstlisting}
                
        or even more simply -
        
        \begin{lstlisting}[gobble=8, language=threepl]
            class(cl, c());
        \end{lstlisting}
        
        Thus a class is called in a similar way to a function but differs in that -
        \blist
        \item   A call to a class only ever returns an immediate value of type
                "class".
        \item   Unlike a function call, a class call can have output arguments.
        \blend        
        
        Variables within the above class instance {\bf cl} can be accessed as
        in the following examples -
        
        \begin{lstlisting}[gobble=8, language=threepl]
            j = cl.i;
            cl.s++;
        \end{lstlisting}
        
        Each class instance also has a variable {\bf this} of type "class"
        defined which is the current class instance. 
        
        Like modules, procedures and functions, variables are re-created on
        each call and hence their values can be different in different instances
        of the class.
        
        It should be noted that parameters created in the call to the class are
        local variables (as is the case with modules, procedures and functions)
        and so they can also be accessed in the same way as the other class
        variables.
        
        There are cases where it would be convenient to share data between class
        instances and this can be done by using a common block -

        \begin{lstlisting}[gobble=8, language=threepl]
            class c () {
                common {
                    value(v, "log");
                }
                int(i, 34);
                static(s, "uint:8");
            }
        \end{lstlisting}
        
        here variable {\bf v} is created once only when the class is called to
        create the first instance of the class. All further instances will share
        the same variable and its value. Common block variables can be both
        evaluated and assigned in any instance of the class and any changes will
        be reflected in all instances of the class. Note that variables are
        only created in the class common block on the first call, they cannot
        be created by later calls.
        
        The common block must be at the top level within the class, i.e. it
        cannot be enclosed within any other control structure.
        
        If common class variables are to be given initial values derived from
        arguments to the first class call then it is probably good practice to
        call the class once initially to create the common class variables and
        employ conditional statements to return immediately having done that. If
        the initialisation call is not intended to create an instance then the
        return value can be discarded by calling the class using a procedure call
        syntax. Later calls to the class will automatically ignore the common
        block. Using key/value argument/parameter assignments makes it easy to
        hide this once-only code, e.g.
        
        \begin{lstlisting}[gobble=8, language=threepl]
            class(cl);              // create variable of type "class"
            
            c(v1=x);                // call to initialise common
            cl = c(i1=5, i2=7));    // call to create a class instance
            
            // class definition
            class c (int({i1, i2}), value(v1, "log")) {
                common {
                    value(v, "log", v1);    // copy v1 to common v
                }
                if (matched(v1)
                    return; // is initial call - ignore the rest.
                    
                int(i, 34);
                static(s, "uint:8");
                .
                i = i1 + i2;
                .
                .
            }
        \end{lstlisting}
        The above could also have instead -
        \begin{lstlisting}[gobble=8, language=threepl]
                if (!matched(i1) && !matched(i2))
                    return;
        \end{lstlisting}
                
        A class can contain definitions of modules, procedures, functions and
        classes and these can be called as in the following example -

        \begin{lstlisting}[gobble=8, language=threepl]
            class(cl, c());
            int(j);
            static(ss, "uint:8");
            
            cl.p(j <- 4);   // assigns 4 * cl.i to j
            ss <- cl.f();   // FPGA code assigning cl.s + 5 to ss
            
            class c () {
                int(i, 34);
                static(s, "uint:8");
                
                procedure p (int(out) <- int(in)) {
                    out = in * i;
                }
                
                function f () {
                    return(s + 5);
                }
            }
        \end{lstlisting}
        
        Here the first assignment statement is immediate, calling class
        procedure {\bf p()}. Parameter {\bf in}, whose argument is 4, is
        mutiplied by class variable {\bf i} and the result is assigned via parameter
        {\bf out} to {\bf j}.
        
        The second assignment statement is target (FPGA) code. The class
        function {\bf f()} is called and returns the sum of the static 8-bit
        unsigned integer class variable {\bf s} and constant 5 and the returned
        value is assigned to static variable {\bf ss}.
        
        A class cannot contain in-line target statements in its code body,
        either directly or via a procedure call.
        
        Modules, procedures and functions can have their own local variables as
        usual but also have read and write access to variables in the class in
        which they are defined (as in the example above). They also can call
        other modules, procedures, functions and classes defined within the
        class.
        
        A module, procedure, function or class defined within a class may
        be preceded by the keyword {\bf private}, e.g.
        \begin{lstlisting}[gobble=8, language=threepl]
            class c () {
                .
                .
                .
                
                private
                procedure p (int(out) <- int(in)) {
                   .
                   .
                }
                
                function f () {
                    .
                    .
                }
            }
        \end{lstlisting}
        which means that the module etc can only be called from within the class,
        not via class instance. For example in -
        
        \begin{lstlisting}[gobble=8, language=threepl]
            class(cl, c());
            .
            cl.p();
            i = cl.f();
            
            class c () {
                .
                .
                
                private
                procedure p (int(out) <- int(in)) {
                    out = in * i;
                }
                
                function f () {
                    return(s + 5);
                }
            }
        \end{lstlisting}
        the call {\bf cl.p()} will fail whereas {\bf i = cl.f()} will succeed.

